\chapter{BX-16: An Amiga-Compatible Machine with SVM-16}
\label{ch:bx16}

\section{The board-level picture}

BX-16 can replace an Amiga in the same practical sense that BX-8 can replace an MSX. The bus is sixteen bits wide. Several masters share it: the hart, a blitter-class block mover, video fetch, and audio DMA. Chip memory and a larger local memory sit on that bus. The display timing can follow a programmed list of fetches, in the role a copper list fills on the original machine. The hart that runs the list is RV32I with a sixteen-bit data path.

Integer registers remain thirty-two bits. A load word takes two halfword beats. Instruction fetch is again allowed its own port so that video DMA does not stall the decode of every parcel.

\section{SVM-16}

\begin{spec}
ELEN = 16, VLEN = 128, thirty-two registers.\\
Legal SEW = 8 or 16.\\
LMUL may be greater than 1 only for a widening destination.\\
Required, beyond SVM-8: strides, widening multiply,\\
narrowing shift, saturating add and subtract,\\
integer add reduction, slide by one, \texttt{vsetivli}.\\
No floating point. Xsvm optional.
\end{spec}

A register holds sixteen elements of eight bits or eight elements of sixteen bits. That is enough for a short span of a scanline or a frame of samples. Longer spans strip-mine, and the bus arbiter inserts the video and audio masters between strips. Write the arbiter so that a maximum latency for the audio master is enforced. A vector copy that can lock out audio is not an Amiga-compatible bus, however correct its RVV tests are.

\section{Widening, not a private fixed-point format}

A sixteen-bit multiply that keeps the full product uses \texttt{vwmul.vv}. The destination elements are thirty-two bits wide, in a register group. A following narrowing shift brings a saturated sixteen-bit pixel back to SEW=16. Software that wants a Q16.16 convention shifts by sixteen. The lane does not test a mode bit named Q16.16.

This is the operation a blitter-class scaler and a small transform both need. Implementing it as a special coprocessor instruction would make the kernel unreadable on BX-32. Implementing it as \texttt{vwmul} means the same source runs on Zve32f.

\section{Xsvm}

When the elaboration parameter \texttt{xsvm} is true, two custom-0 instructions are decoded.

The matrix instruction reads four vector registers starting at a fixed base, \texttt{v8} through \texttt{v11}, as the rows of a four-by-four matrix, multiplies them by \texttt{vs2}, and writes \texttt{vd}. Using \texttt{v8} rather than \texttt{v0} keeps the mask register free. The cross-product instruction reads two registers and writes one, three elements, at the current SEW, with \texttt{vl} forced to a minimum of three or the instruction is illegal.

Both instructions are undisturbed outside \texttt{vl}. Both are absent from the decoder when \texttt{xsvm} is false, and the parcels trap. The exercises that must pass on every BX-16 image do not use them.

\section{Opcode map for SVM-16}

Everything legal on SVM-8 remains legal. The additional classes required by SVM-16 are shown with the same rigid two-line matrix. SEW may be 8 or 16. The complete cross-profile contract and full opcode summary appear after the four machines, in Chapter~\ref{ch:bsv}.

\subsection{Configuration}

Identical to SVM-8. \texttt{vsetivli} is now required.

\subsection{Unit-stride load/store (EEW = 8 and 16)}

Same layout as SVM-8. Width field:
\begin{itemize}[leftmargin=1.5em]
  \item 000 = 8-bit (\texttt{vle8.v}/\texttt{vse8.v})
  \item 101 = 16-bit (\texttt{vle16.v}/\texttt{vse16.v})
\end{itemize}

\subsection{Constant-stride load --- \texttt{vlse16.v}}

\begin{center}
\ttfamily\footnotesize
\setlength{\tabcolsep}{0.3em}
\begin{tabular}{*{16}{w{c}{1.05em}}}
31 & 30 & 29 & 28 & 27 & 26 & 25 & 24 & 23 & 22 & 21 & 20 & 19 & 18 & 17 & 16 \\
\hline
\multicolumn{3}{|c|}{N} &
\multicolumn{1}{c|}{E} &
\multicolumn{2}{c|}{P} &
\multicolumn{1}{c|}{M} &
\multicolumn{5}{c|}{S} &
\multicolumn{4}{c|}{B} \\
\hline
\end{tabular}
\vspace{0.35ex}
\begin{tabular}{*{16}{w{c}{1.05em}}}
15 & 14 & 13 & 12 & 11 & 10 & 9 & 8 & 7 & 6 & 5 & 4 & 3 & 2 & 1 & 0 \\
\hline
\multicolumn{3}{|c|}{W} &
\multicolumn{5}{c|}{D} &
\multicolumn{7}{c|}{O} \\
\hline
\end{tabular}
\end{center}

\vspace{1.4ex}
\begin{center}
\begin{tabular}{@{}ll@{}}
N & = nf \\
E & = mew \\
P & = mop (10 = strided) \\
M & = vm \\
S & = rs2 (stride) \\
B & = rs1 (base) \\
W & = width (101 = 16-bit) \\
D & = vd \\
O & = 0000111 \\
\end{tabular}
\end{center}

Store form uses major opcode 0100111.

\subsection{Integer arithmetic (SEW = 8, 16)}

Same OP-V layout as SVM-8. Additional instructions:

\begin{center}
\begin{tabular}{@{}llc@{}}
\toprule
Instruction & funct6 & notes \\
\midrule
\texttt{vminu}/\texttt{vmin} & 000100 / 000101 & vv / vx \\
\texttt{vmaxu}/\texttt{vmax} & 000110 / 000111 & vv / vx \\
\texttt{vsaddu}/\texttt{vsadd} & 100000 / 100001 & saturating \\
\texttt{vssubu}/\texttt{vssub} & 100010 / 100011 & saturating \\
\bottomrule
\end{tabular}
\end{center}

\subsection{Widening multiply and narrowing}

\begin{center}
\ttfamily\footnotesize
\setlength{\tabcolsep}{0.3em}
\begin{tabular}{*{16}{w{c}{1.05em}}}
31 & 30 & 29 & 28 & 27 & 26 & 25 & 24 & 23 & 22 & 21 & 20 & 19 & 18 & 17 & 16 \\
\hline
\multicolumn{6}{|c|}{F} &
\multicolumn{1}{c|}{M} &
\multicolumn{5}{c|}{S} &
\multicolumn{4}{c|}{T} \\
\hline
\end{tabular}
\vspace{0.35ex}
\begin{tabular}{*{16}{w{c}{1.05em}}}
15 & 14 & 13 & 12 & 11 & 10 & 9 & 8 & 7 & 6 & 5 & 4 & 3 & 2 & 1 & 0 \\
\hline
\multicolumn{1}{|c|}{T} &
\multicolumn{3}{c|}{C} &
\multicolumn{5}{c|}{D} &
\multicolumn{7}{c|}{O} \\
\hline
\end{tabular}
\end{center}

\vspace{1.4ex}
\begin{center}
\begin{tabular}{@{}ll@{}}
F & = funct6 \\
M & = vm \\
S & = vs2 \\
T & = vs1 / rs1 \\
C & = funct3 (OPIVV / OPIVX) \\
D & = vd (group for widening) \\
O & = 1010111 \\
\end{tabular}
\end{center}

\begin{center}
\begin{tabular}{@{}llc@{}}
\toprule
Instruction & funct6 & width rule \\
\midrule
\texttt{vwmul}   & 111011 & SEW $\to$ 2$\times$SEW \\
\texttt{vwmulu}  & 111000 & unsigned \\
\texttt{vwmulsu} & 111010 & mixed \\
\texttt{vnsrl}   & 101100 & narrowing shift \\
\texttt{vnsra}   & 101101 & narrowing arithmetic \\
\texttt{vnclipu}/\texttt{vnclip} & 101110 / 101111 & narrowing clip \\
\bottomrule
\end{tabular}
\end{center}

\subsection{Reduction and slide}

\begin{center}
\begin{tabular}{@{}llc@{}}
\toprule
Instruction & funct6 / notes \\
\midrule
\texttt{vredsum.vs} & 000000 (OPMVV) \\
\texttt{vslide1up.vx} & 001110 (OPIVX) \\
\texttt{vslide1down.vx} & 001111 (OPIVX) \\
\bottomrule
\end{tabular}
\end{center}

\subsection{Optional Xsvm (custom-0)}

\begin{center}
\ttfamily\footnotesize
\setlength{\tabcolsep}{0.3em}
\begin{tabular}{*{16}{w{c}{1.05em}}}
31 & 30 & 29 & 28 & 27 & 26 & 25 & 24 & 23 & 22 & 21 & 20 & 19 & 18 & 17 & 16 \\
\hline
\multicolumn{6}{|c|}{F} &
\multicolumn{1}{c|}{M} &
\multicolumn{5}{c|}{S} &
\multicolumn{4}{c|}{X} \\
\hline
\end{tabular}
\vspace{0.35ex}
\begin{tabular}{*{16}{w{c}{1.05em}}}
15 & 14 & 13 & 12 & 11 & 10 & 9 & 8 & 7 & 6 & 5 & 4 & 3 & 2 & 1 & 0 \\
\hline
\multicolumn{1}{|c|}{X} &
\multicolumn{3}{c|}{C} &
\multicolumn{5}{c|}{D} &
\multicolumn{7}{c|}{O} \\
\hline
\end{tabular}
\end{center}

\vspace{1.4ex}
\begin{center}
\begin{tabular}{@{}ll@{}}
F & = custom funct6 \\
M & = vm \\
S & = vs2 \\
X & = unused / fixed \\
C & = custom funct3 \\
D & = vd \\
O & = 0001011 (custom-0) \\
\end{tabular}
\end{center}

Matrix and cross-product are distinguished by the custom funct6 / funct3 values documented in the profile card. They are illegal when \texttt{xsvm = False}.

\section{Assembler kernels that must run}

In addition to every kernel of Chapter~\ref{ch:bx8}, the following must pass on every BX-16 image (Xsvm kernels are excluded from the travelling suite).

\subsection{sew16 --- halfword fill}
\begin{lstlisting}
# a0 = n (elements), a1 = fill halfword, a2 = dst
sew16:
    beqz    a0, 2f
1:  vsetvli t0, a0, e16, m1, tu, mu
    vmv.v.x v1, a1
    vse16.v v1, (a2)
    slli    t1, t0, 1
    add     a2, a2, t1
    sub     a0, a0, t0
    bnez    a0, 1b
2:  ret
\end{lstlisting}

\subsection{vwmul16 --- widening multiply then narrow with saturation}
\begin{lstlisting}
# a0 = n, a1 = src0 (SEW16), a2 = src1, a3 = dst (SEW16)
vwmul16:
    beqz    a0, 2f
1:  vsetvli t0, a0, e16, m1, tu, mu
    vle16.v v2, (a1)
    vle16.v v4, (a2)
    vwmul.vv v8, v2, v4
    vsetvli t0, a0, e16, m1, tu, mu
    vnclip.wi v1, v8, 0
    vse16.v v1, (a3)
    slli    t1, t0, 1
    add     a1, a1, t1
    add     a2, a2, t1
    add     a3, a3, t1
    sub     a0, a0, t0
    bnez    a0, 1b
2:  ret
\end{lstlisting}

\subsection{stride16 --- constant-stride gather of halfwords}
\begin{lstlisting}
# a0 = n, a1 = src base, a2 = stride (bytes), a3 = dst
stride16:
    beqz    a0, 2f
1:  vsetvli t0, a0, e16, m1, tu, mu
    vlse16.v v1, (a1), a2
    vse16.v  v1, (a3)
    slli    t1, t0, 1
    add     a3, a3, t1
    sub     a0, a0, t0
    bnez    a0, 1b
2:  ret
\end{lstlisting}

\subsection{redsum16 --- integer reduction}
\begin{lstlisting}
# a0 = n, a1 = src (SEW16), a2 = scalar accumulator out
redsum16:
    vsetvli t0, a0, e16, m1, tu, mu
    vle16.v v1, (a1)
    vmv.v.i v2, 0
    vredsum.vs v2, v1, v2
    vmv.x.s a2, v2
    ret
\end{lstlisting}

\section{One hart}

Scalar parcels and vector parcels share decode, the scoreboard, and writeback. The blitter, if it is a bus master, is not a second hart and does not fetch OP-V. It is a peripheral with a pointer, a count, and a small arithmetic step, programmed by stores to its registers. New software may ignore it and use SVM-16 instead. Old software that you have translated may still program it. Keeping both is how the machine remains Amiga-compatible while the new kernels are ordinary RVV.

\section{Elaboration}

\begin{lstlisting}
Profile profile = SVM16;
Integer elen = 16;
Integer vlen = 128;
Integer lanes = 4;
Bool xsvm = False; // True only on a geometry image
\end{lstlisting}

Four lanes of sixteen bits retire four elements per cycle. A full SEW=16 register then takes two arithmetic cycles, plus the read and write of the file. Publish that number. Programmers of audio mixing will plan their strip length around it.
